% --- Archon physics formalization source begin ---
% archon:physics
% archon:covers IPhO2026Problems/problem_IPhO_2026_1_C_1.lean
% archon:source-report reports/ipho_2026/problem_IPhO_2026_1_C_1.source.json
% archon:problem-id IPhO_2026_1
% archon:part-id C.1

\chapter{Physics problem IPhO\_2026\_1\_C\_1}
\label{ch:IPhO2026Problems_problem_IPhO_2026_1_C_1}

\paragraph{Problem source.}
A photon of angular frequency omega is absorbed by an ozone molecule O3 at rest,
dissociating it into O2 and O.  Let U\_i and U\_f be the ground-state energies of
O3 and O2 and define Delta U = U\_f - U\_i.  The outgoing O2 momentum makes angle
theta with the incident photon.  Treat the oxygen fragments classically and
non-relativistically, take the mass of an oxygen atom to be m, and use photon
momentum p\_gamma = E\_gamma/c = hbar*omega/c.

Current subquestion:
Determine the minimum angular frequency omega\_min required for dissociation at outgoing O2 angle theta, in terms of hbar, c, theta, Delta U, and m.

\paragraph{Current subquestion.}
Determine the minimum angular frequency omega\_min required for dissociation at outgoing O2 angle theta, in terms of hbar, c, theta, Delta U, and m.

\paragraph{Recorded answer/context.}
The source-report transcription omits a factor \(2\) under the square root.
The conservation laws instead give, for \(A=2\sin^2\theta+1\) and
\(\theta\leq\pi/2\),
\[
\omega_{\min}=\frac{3mc^2}{\hbar A}
  \left(1-\sqrt{1-\frac{2A\Delta U}{3mc^2}}\right).
\]
For \(\theta\geq\pi/2\), use the constrained value at \(\theta=\pi/2\):
\[
\omega_{\min}=\frac{mc^2}{\hbar}
  \left(1-\sqrt{1-\frac{2\Delta U}{mc^2}}\right).
\]

\paragraph{Figure/image path.}
/root/proposal\_for\_physic/science-mango/ipho\_2026\_source/image/T1\_page-3.png

\paragraph{Formalization target.}
Redraft the existing compiling scaffold so its conclusion contains the corrected
factor \(2\), while preserving the dimensioned masses, energy gap, action,
frequency, momenta, Figure 1c angle, conservation laws, feasibility predicate,
and minimum-frequency meaning.  Leave proof bodies as `by sorry`.

\begin{definition}[MassQuantity]
  \label{decl:physics:IPhO_2026_1_C_1:MassQuantity}
  \lean{IPhO2026Problems.IPhO2026_1_C_1.MassQuantity}
  The dimensionful mass of one oxygen atom.
\end{definition}
\begin{proof}
  This is the defining declaration; unfolding it gives the stated typed data or relation.
\end{proof}

\begin{definition}[ActionQuantity]
  \label{decl:physics:IPhO_2026_1_C_1:ActionQuantity}
  \lean{IPhO2026Problems.IPhO2026_1_C_1.ActionQuantity}
  A dimensionful action, used for the reduced Planck constant.
\end{definition}
\begin{proof}
  This is the defining declaration; unfolding it gives the stated typed data or relation.
\end{proof}

\begin{definition}[AngularFrequencyQuantity]
  \label{decl:physics:IPhO_2026_1_C_1:AngularFrequencyQuantity}
  \lean{IPhO2026Problems.IPhO2026_1_C_1.AngularFrequencyQuantity}
  A dimensionful angular frequency.
\end{definition}
\begin{proof}
  This is the defining declaration; unfolding it gives the stated typed data or relation.
\end{proof}

\begin{definition}[MomentumQuantity2]
  \label{decl:physics:IPhO_2026_1_C_1:MomentumQuantity2}
  \lean{IPhO2026Problems.IPhO2026_1_C_1.MomentumQuantity2}
  A dimensionful two-dimensional momentum used for Figure 1c.
\end{definition}
\begin{proof}
  This is the defining declaration; unfolding it gives the stated typed data or relation.
\end{proof}

\begin{definition}[scalarSI]
  \label{decl:physics:IPhO_2026_1_C_1:scalarSI}
  \lean{IPhO2026Problems.IPhO2026_1_C_1.scalarSI}
  The real scalar readout of a dimensionful scalar quantity in SI units.
\end{definition}
\begin{proof}
  This is the defining declaration; unfolding it gives the stated typed data or relation.
\end{proof}

\begin{definition}[speedSI]
  \label{decl:physics:IPhO_2026_1_C_1:speedSI}
  \lean{IPhO2026Problems.IPhO2026_1_C_1.speedSI}
  The real speed readout in metres per second.
\end{definition}
\begin{proof}
  This is the defining declaration; unfolding it gives the stated typed data or relation.
\end{proof}

\begin{definition}[momentumSI]
  \label{decl:physics:IPhO_2026_1_C_1:momentumSI}
  \lean{IPhO2026Problems.IPhO2026_1_C_1.momentumSI}
  \uses{decl:physics:IPhO_2026_1_C_1:MomentumQuantity2}
  The two Cartesian momentum components in SI units.
\end{definition}
\begin{proof}
  This is the defining declaration; unfolding it gives the stated typed data or relation.
\end{proof}

\begin{definition}[dot2]
  \label{decl:physics:IPhO_2026_1_C_1:dot2}
  \lean{IPhO2026Problems.IPhO2026_1_C_1.dot2}
  Euclidean dot product of the two displayed momentum vectors.
\end{definition}
\begin{proof}
  This is the defining declaration; unfolding it gives the stated typed data or relation.
\end{proof}

\begin{definition}[magnitude2]
  \label{decl:physics:IPhO_2026_1_C_1:magnitude2}
  \lean{IPhO2026Problems.IPhO2026_1_C_1.magnitude2}
  Euclidean magnitude of a two-dimensional momentum readout.
\end{definition}
\begin{proof}
  This is the defining declaration; unfolding it gives the stated typed data or relation.
\end{proof}

\begin{definition}[PhotodissociationParameters]
  \label{decl:physics:IPhO_2026_1_C_1:PhotodissociationParameters}
  \lean{IPhO2026Problems.IPhO2026_1_C_1.PhotodissociationParameters}
  \uses{decl:physics:IPhO_2026_1_C_1:MassQuantity, decl:physics:IPhO_2026_1_C_1:ActionQuantity}
  The physical parameters appearing in the problem.
\end{definition}
\begin{proof}
  This is the defining declaration; unfolding it gives the stated typed data or relation.
\end{proof}

\begin{definition}[reducedPlanckConstantSI]
  \label{decl:physics:IPhO_2026_1_C_1:reducedPlanckConstantSI}
  \lean{IPhO2026Problems.IPhO2026_1_C_1.reducedPlanckConstantSI}
  \uses{decl:physics:IPhO_2026_1_C_1:scalarSI, decl:physics:IPhO_2026_1_C_1:PhotodissociationParameters}
  The SI readout of “ℏ”.
\end{definition}
\begin{proof}
  This is the defining declaration; unfolding it gives the stated typed data or relation.
\end{proof}

\begin{definition}[lightSpeedSI]
  \label{decl:physics:IPhO_2026_1_C_1:lightSpeedSI}
  \lean{IPhO2026Problems.IPhO2026_1_C_1.lightSpeedSI}
  \uses{decl:physics:IPhO_2026_1_C_1:speedSI, decl:physics:IPhO_2026_1_C_1:PhotodissociationParameters}
  The SI readout of the speed of light.
\end{definition}
\begin{proof}
  This is the defining declaration; unfolding it gives the stated typed data or relation.
\end{proof}

\begin{definition}[oxygenAtomMassSI]
  \label{decl:physics:IPhO_2026_1_C_1:oxygenAtomMassSI}
  \lean{IPhO2026Problems.IPhO2026_1_C_1.oxygenAtomMassSI}
  \uses{decl:physics:IPhO_2026_1_C_1:scalarSI, decl:physics:IPhO_2026_1_C_1:PhotodissociationParameters}
  The SI readout of the mass “m” of one oxygen atom.
\end{definition}
\begin{proof}
  This is the defining declaration; unfolding it gives the stated typed data or relation.
\end{proof}

\begin{definition}[energyDifferenceSI]
  \label{decl:physics:IPhO_2026_1_C_1:energyDifferenceSI}
  \lean{IPhO2026Problems.IPhO2026_1_C_1.energyDifferenceSI}
  \uses{decl:physics:IPhO_2026_1_C_1:scalarSI, decl:physics:IPhO_2026_1_C_1:PhotodissociationParameters}
  The energy gap “ΔU = U\_f - U\_i”, in joules.
\end{definition}
\begin{proof}
  This is the defining declaration; unfolding it gives the stated typed data or relation.
\end{proof}

\begin{definition}[ValidPhotodissociationParameters]
  \label{decl:physics:IPhO_2026_1_C_1:ValidPhotodissociationParameters}
  \lean{IPhO2026Problems.IPhO2026_1_C_1.ValidPhotodissociationParameters}
  \uses{decl:physics:IPhO_2026_1_C_1:PhotodissociationParameters, decl:physics:IPhO_2026_1_C_1:reducedPlanckConstantSI, decl:physics:IPhO_2026_1_C_1:lightSpeedSI, decl:physics:IPhO_2026_1_C_1:oxygenAtomMassSI, decl:physics:IPhO_2026_1_C_1:energyDifferenceSI, decl:physics:IPhO_2026_1_C_1:DissociationAt}
  Physical parameter and angle conditions.
\end{definition}
\begin{proof}
  This is the defining declaration; unfolding it gives the stated typed data or relation.
\end{proof}

\begin{definition}[DissociationAt]
  \label{decl:physics:IPhO_2026_1_C_1:DissociationAt}
  \lean{IPhO2026Problems.IPhO2026_1_C_1.DissociationAt}
  \uses{decl:physics:IPhO_2026_1_C_1:MomentumQuantity2, decl:physics:IPhO_2026_1_C_1:momentumSI, decl:physics:IPhO_2026_1_C_1:dot2, decl:physics:IPhO_2026_1_C_1:magnitude2, decl:physics:IPhO_2026_1_C_1:PhotodissociationParameters, decl:physics:IPhO_2026_1_C_1:reducedPlanckConstantSI, decl:physics:IPhO_2026_1_C_1:lightSpeedSI, decl:physics:IPhO_2026_1_C_1:oxygenAtomMassSI, decl:physics:IPhO_2026_1_C_1:energyDifferenceSI}
  “DissociationAt p θ ω” states the governing laws for a photon of angular frequency “ω”, measured in radians per second, producing the two fragments at the Figure 1c angle “θ”.
\end{definition}
\begin{proof}
  This is the defining declaration; unfolding it gives the stated typed data or relation.
\end{proof}

\begin{definition}[IsMinimumDissociationFrequency]
  \label{decl:physics:IPhO_2026_1_C_1:IsMinimumDissociationFrequency}
  \lean{IPhO2026Problems.IPhO2026_1_C_1.IsMinimumDissociationFrequency}
  \uses{decl:physics:IPhO_2026_1_C_1:AngularFrequencyQuantity, decl:physics:IPhO_2026_1_C_1:scalarSI, decl:physics:IPhO_2026_1_C_1:PhotodissociationParameters, decl:physics:IPhO_2026_1_C_1:DissociationAt}
  The proposed angular frequency is feasible and no smaller feasible nonnegative frequency exists at the same outgoing “O₂” angle.
\end{definition}
\begin{proof}
  This is the defining declaration; unfolding it gives the stated typed data or relation.
\end{proof}

\begin{theorem}[Physics formalization target]
\label{thm:physics:IPhO_2026_1_C_1:target}
\lean{IPhO2026Problems.IPhO2026_1_C_1.minimumAngularFrequency_eq}
\uses{decl:physics:IPhO_2026_1_C_1:MassQuantity, decl:physics:IPhO_2026_1_C_1:ActionQuantity, decl:physics:IPhO_2026_1_C_1:AngularFrequencyQuantity, decl:physics:IPhO_2026_1_C_1:MomentumQuantity2, decl:physics:IPhO_2026_1_C_1:scalarSI, decl:physics:IPhO_2026_1_C_1:speedSI, decl:physics:IPhO_2026_1_C_1:momentumSI, decl:physics:IPhO_2026_1_C_1:dot2, decl:physics:IPhO_2026_1_C_1:magnitude2, decl:physics:IPhO_2026_1_C_1:PhotodissociationParameters, decl:physics:IPhO_2026_1_C_1:reducedPlanckConstantSI, decl:physics:IPhO_2026_1_C_1:lightSpeedSI, decl:physics:IPhO_2026_1_C_1:oxygenAtomMassSI, decl:physics:IPhO_2026_1_C_1:energyDifferenceSI, decl:physics:IPhO_2026_1_C_1:ValidPhotodissociationParameters, decl:physics:IPhO_2026_1_C_1:DissociationAt, decl:physics:IPhO_2026_1_C_1:IsMinimumDissociationFrequency}
For valid dimensioned photodissociation data, the least feasible photon
frequency equals the corrected forward-angle expression above when
\(\theta\leq\pi/2\), and the boundary expression when
\(\theta\geq\pi/2\).
\end{theorem}
\begin{proof}
Write \(q=\hbar\omega/c\) for the incident photon momentum magnitude and
\(p\geq0\) for the outgoing molecular momentum magnitude.  Momentum
conservation makes the atomic momentum the vector difference, so the total
fragment kinetic energy is
\[
\frac{3p^2-4pq\cos\theta+2q^2}{4m}.
\]
If \(\cos\theta\geq0\), this quadratic is minimized at
\(p=2q\cos\theta/3\), giving \(q^2(2\sin^2\theta+1)/(6m)\).
Substitute \(q=\hbar\omega/c\) into energy conservation and take the smaller
nonnegative root of the resulting quadratic.  If \(\cos\theta\leq0\), the
constraint \(p\geq0\) puts the minimum at \(p=0\); solving the corresponding
quadratic gives the boundary formula.
\end{proof}
% --- Archon physics formalization source end ---
